\subsection{异常定位}

通过将 EDR 受激序列 \(X^{E}\) 输入基线行为模型中识别 EDR 在恶意刺激下产生的异常响应。具体而言，我们采用多尺度窗口并基于负对数似然分数进行异常定位。
设窗口长度集合为：
\begin{equation}
    \mathcal{W} = \{256, 512, 1024\}.
\end{equation}

对于每个时间位置 \(t\)，分别计算不同窗口长度下的平均负对数似然，并聚合得到多尺度异常分数：
\begin{equation}
    S(t)
    =
    \max_{w \in \mathcal{W}}
    \mathrm{NLL}_{w}(t).
\end{equation}

由于不同 EDR 产品的异常行为强度与模型 NLL 尺度存在差异，需要为每个目标产品 \(p\) 自适应设置异常阈值。因此，我们首先在仅运行目标 EDR 的基线序列 \(X^{B}\) 上计算多尺度异常分数：
\begin{equation}
    S^{B}(t)
    =
    \max_{w \in \mathcal{W}}
    \mathrm{NLL}^{B}_{w}(t).
\end{equation}

随后，根据基线异常分数的经验分布确定对应产品阈值：
\begin{equation}
    \tau_p
    =
    Q_q\left(\{S^{B}(t)\}_{t}\right),
\end{equation}
其中，\(Q_q(\cdot)\) 表示经验分布的第 \(q\) 分位数。采用 \(q=0.995\)，即将正常基线运行中最高 \(0.5\%\) 的窗口作为可容忍的背景波动上界。该设置使阈值能够随不同 EDR 产品及其运行配置自适应变化，并在基线条件下控制异常窗口的误报比例。

对于受激序列 \(X^{E}\)，若某个时间位置 \(t\) 的多尺度异常分数满足：
\begin{equation}
    S^{E}(t) > \tau_p,
\end{equation}
则认为该位置对应的 EDR 行为窗口显著偏离产品 \(p\) 的基线行为模型。由此得到异常窗口集合：
\begin{equation}
    \mathcal{A}
    =
    \{a_1, a_2, \ldots, a_m\}.
\end{equation}
